\SourceLeaf{321}
\[
q=-\cos\chi\mathord{.}x-\tfrac12\operatorname{tang}Q\mathord{.}\sin\chi^2xx\quad .\ .\ .
\]
(deren leichte Ableitung hier weggelassen werden kann), und das Resultat mit der Reihe
\[
m=1+\mu x^3+\mu'x^4\quad .\ .\ .
\]
zusammenhält. Für unsern gegenwärtigen Zweck ist jedoch mehr nicht nöthig, als nachzuweisen, dass die gesuchte Länge des geodaetischen Bogens vom \(Ah\) nicht mehr als um eine Grösse fünfter Ordnung abweicht. Da nun ersichtlich \(h^5+10h^3\delta\delta+5h\delta^4\) eine solche Grösse ist, so braucht der entwickelte Werth von \(\mu'\) nicht hiehergesetzt zu werden. Für \(\mu\) aber ergibt sich der Werth
\[
\mu=\frac{2ee\cos P\mathord{.}\sin P}{3\cos\varphi\cos\theta}\mathord{.}\cos\chi^3
\]
und da \(\delta\cos\chi\) nach Art. 15 eine Grösse zweiter Ordnung ist, so wird offenbar auch \((h^3\delta+h\delta^3)\mu\) eine Grösse fünfter Ordnung.

Wir haben demnach, da \(h\) dasselbe bedeutet, was jetzt mit \(s\) bezeichnet ist, bis auf die fünfte Ordnung ausschliesslich genau
\[
s=\frac rA\quad .\ .\ .\ .\ .\ .\ .\ .\ .\ .\ .\ .\ .\ .\ .\quad(8)
\]

\indent Endlich, damit alles für die weitere Entwicklung erforderliche hier beisammen sei, setze ich noch folgende schon in der ersten Abhandlung (Art. 4, 6 und 3) gebrauchte strenge richtige Gleichungen hieher:
\[
\begin{aligned}
A&=\frac{a\cos\varphi}{\cos\theta^2}\quad .\ .\ .\ .\ .\ .\ .\ .\ .\ .\ .\ .\ .\ .\quad(9)\\
\cos Q&=\frac{\cos\theta\cos P}{\alpha\cos\varphi}\quad .\ .\ .\ .\ .\ .\ .\ .\ .\ .\ .\ .\ .\quad(10)\\
\sin Q&=\frac{\sin P}{\alpha}
\end{aligned}
\]
und die aus der Verbindung dieser beiden hervorgehende
\[
\operatorname{tang}Q=\frac{\cos\varphi\operatorname{tang}P}{\cos\theta}\quad .\ .\ .\ .\ .\ .\ .\ .\ .\ .\ .\ .\quad(11)
\]

\Article{23}
Zur Erreichung unsers Zwecks brauchen nun bloss diese Gleichungen gehörig combinirt zu werden.
